\documentclass{handout}

\usepackage{handoutSetup}\usepackage{handoutShortcuts}
 

\pagestyle{empty}

\begin{document}
\begin{verbatimwrite}{\jobname.title}
Gross Saving and Growth in the RCK Model
\end{verbatimwrite}

\handoutHeader
\handoutNameMake

In the neoclassical growth model with labor-augmenting technological
progress at rate $\pGro$, utility function $\util(c) = c^{1-\CRRA}/(1-\CRRA)$,
time preference rate $\DiscRate$ and depreciation rate $\delta$ the
steady-state will be at the point where the growth rate of consumption
is equal to the growth rate of labor-augmenting technological 
progress, $\pGro$,
\begin{equation}\begin{gathered}\begin{aligned}
        \frac{\dot{C}}{C} & =  \CRRA^{-1}(f^{\prime}(k)-\delta-\DiscRate)  \\
         & =  \pGro  
\end{aligned}\end{gathered}\end{equation}
which implies that
\begin{equation}\begin{gathered}\begin{aligned}
        f^{\prime}(k) & =  \CRRA \pGro + \DiscRate + \delta  \\
        \alpha k^{\alpha-1} & =  \CRRA \pGro + \DiscRate + \delta  \\
    k & =  \left[\frac{\CRRA \pGro + \DiscRate + \delta}{\alpha}\right]^{\frac{1}{\alpha-1}}. 
\end{aligned}\end{gathered}\end{equation}

The aggregate gross saving rate is defined as
\begin{equation}\begin{gathered}\begin{aligned}
     s & =     \frac{y-c}{y} \\
         & =  \frac{k^{\alpha}-c}{k^{\alpha}} \\
         & =  1-c/k^{\alpha}. \label{eq:grossavrate}
\end{aligned}\end{gathered}\end{equation}

In steady-state by definition
\begin{equation}\begin{gathered}\begin{aligned}
        \dot{k} & =  0 \label{eq:kdotOk} 
\end{aligned}\end{gathered}\end{equation}
but from the capital accumulation equation we know that
\begin{equation}\begin{gathered}\begin{aligned}
        \dot{k} & =  k^{\alpha}-(\delta+\pGro)k - c \label{eq:dotk}
\end{aligned}\end{gathered}\end{equation}
so in steady-state
\begin{equation}\begin{gathered}\begin{aligned}
        c & =  k^{\alpha}-(\delta+\pGro)k.
\end{aligned}\end{gathered}\end{equation}

This can be substituted into (\ref{eq:grossavrate}) to obtain
\begin{equation}\begin{gathered}\begin{aligned}
        s & =  1-\frac{k^{\alpha}-(\delta+\pGro)k}{k^{\alpha}}  \\
         & =  (\delta+\pGro)k^{1-\alpha}
\end{aligned}\end{gathered}\end{equation}
and the expression for the steady-state level of capital per capita
can be substituted in to yield
\begin{equation}\begin{gathered}\begin{aligned}
        s & =  (\delta+\pGro)\left(\frac{\CRRA \pGro+ \DiscRate + \delta}{\alpha}\right)^{-1}  
\\ & =  \left(\frac{\alpha(\pGro+\delta)}{\delta+\DiscRate+\CRRA \pGro}\right)     
\end{aligned}\end{gathered}\end{equation}
\clearpage

The derivative of this expression with respect to $\pGro$ is
\begin{equation}\begin{gathered}\begin{aligned}
        \frac{ds}{d\pGro} & =  \left(\frac{\alpha(\delta + \DiscRate + \CRRA \pGro)-\alpha 
(\pGro+\delta) \CRRA}{(\delta+\DiscRate+\CRRA \pGro)^{2}}\right)
\\ & =  \left(\frac{\alpha(\delta(1-\CRRA) + \DiscRate)}{(\delta+\DiscRate+\CRRA 
\pGro)^{2}}\right)
.
\end{aligned}\end{gathered}\end{equation}

This will be positive if its numerator is positive, i.e.\ if
\begin{equation}\begin{gathered}\begin{aligned}
         \CRRA \delta & <  \DiscRate + \delta
\\  \CRRA & <  1+\DiscRate/\delta  .
\end{aligned}\end{gathered}\end{equation}

A typical assumption is $\DiscRate = .04$ and $\delta = .08$, implying that 
the steady-state relationship between saving and growth in the neoclassical
model is positive only if the coefficient of relative risk aversion $\CRRA$
is less than 1.5.  Typically we assume values of $\CRRA$ in the range from
2 to 5, so the model leads us to expect a negative relationship between saving
and growth.

\end{document}
